% Options for packages loaded elsewhere
\PassOptionsToPackage{unicode}{hyperref}
\PassOptionsToPackage{hyphens}{url}
\PassOptionsToPackage{dvipsnames,svgnames,x11names}{xcolor}
\documentclass[
  11pt,
]{article}
\usepackage{xcolor}
\usepackage[margin=25mm]{geometry}
\usepackage{amsmath,amssymb}
\setcounter{secnumdepth}{-\maxdimen} % remove section numbering
\usepackage{iftex}
\ifPDFTeX
  \usepackage[T1]{fontenc}
  \usepackage[utf8]{inputenc}
  \usepackage{textcomp} % provide euro and other symbols
\else % if luatex or xetex
  \usepackage{unicode-math} % this also loads fontspec
  \defaultfontfeatures{Scale=MatchLowercase}
  \defaultfontfeatures[\rmfamily]{Ligatures=TeX,Scale=1}
\fi
\usepackage{lmodern}
\ifPDFTeX\else
  % xetex/luatex font selection
\fi
% Use upquote if available, for straight quotes in verbatim environments
\IfFileExists{upquote.sty}{\usepackage{upquote}}{}
\IfFileExists{microtype.sty}{% use microtype if available
  \usepackage[]{microtype}
  \UseMicrotypeSet[protrusion]{basicmath} % disable protrusion for tt fonts
}{}
\makeatletter
\@ifundefined{KOMAClassName}{% if non-KOMA class
  \IfFileExists{parskip.sty}{%
    \usepackage{parskip}
  }{% else
    \setlength{\parindent}{0pt}
    \setlength{\parskip}{6pt plus 2pt minus 1pt}}
}{% if KOMA class
  \KOMAoptions{parskip=half}}
\makeatother
\usepackage{longtable,booktabs,array}
\newcounter{none} % for unnumbered tables
\usepackage{calc} % for calculating minipage widths
% Correct order of tables after \paragraph or \subparagraph
\usepackage{etoolbox}
\makeatletter
\patchcmd\longtable{\par}{\if@noskipsec\mbox{}\fi\par}{}{}
\makeatother
% Allow footnotes in longtable head/foot
\IfFileExists{footnotehyper.sty}{\usepackage{footnotehyper}}{\usepackage{footnote}}
\makesavenoteenv{longtable}
\setlength{\emergencystretch}{3em} % prevent overfull lines
\providecommand{\tightlist}{%
  \setlength{\itemsep}{0pt}\setlength{\parskip}{0pt}}
\usepackage{bookmark}
\IfFileExists{xurl.sty}{\usepackage{xurl}}{} % add URL line breaks if available
\urlstyle{same}
\hypersetup{
  pdftitle={Privacy is Value · V6: The Mage Reading},
  pdfauthor={privacymage},
  colorlinks=true,
  linkcolor={black},
  filecolor={Maroon},
  citecolor={Blue},
  urlcolor={blue},
  pdfcreator={LaTeX via pandoc}}

\title{Privacy is Value · V6: The Mage Reading}
\usepackage{etoolbox}
\makeatletter
\providecommand{\subtitle}[1]{% add subtitle to \maketitle
  \apptocmd{\@title}{\par {\large #1 \par}}{}{}
}
\makeatother
\subtitle{Companion guide to the Privacy Value Model · context,
narrative, standards, economics · a volume of the Privacy is Value book}
\author{privacymage}
\date{2026-06-10}

\begin{document}
\maketitle

\section{Privacy Value Model V6: Companion
Guide}\label{privacy-value-model-v6-companion-guide}

\subsection{How to Read These
Documents}\label{how-to-read-these-documents}

Three documents carry V6, one model. The \textbf{formal specification}
(\texttt{privacy\_value\_v6.md}) is the authority: it mirrors the V5.4
skeleton section for section, marks every section CARRIED, REVISED, or
NEW, and cites every conjecture from the unified register. The
\textbf{compressed specification} (\texttt{pvm\_v6\_compressed.md}) is
the Swordsman reading: equations only, five pages, for agents and
reviewers who need the shape before the proofs. This guide is the Mage
reading: what the mathematics means, where it came from, and what it
asks of you. Under the unified-V6 labeling decision, the research paper
and the whitepaper carry the same V6 label; this guide tells you when to
reach for each.

\subsection{1. The Two Strands of V6}\label{the-two-strands-of-v6}

V6 was planned before it was discovered, and the plan came first.

\textbf{The gathering turn.} From the earliest V6 research notes, V6 was
identified as the version that opens the model outward. V5 answered WHAT
the architecture is: two agents, a boundary, a gate, a lattice. V6 asks
WHO. It shifts the document's address into the second person, opens the
registers to the City of Mages and to kindred builders, and turns the
equation from a specification into an invitation, so that its terms can
be filled with data instead of estimates. The betweenness centrality of
the gap got a computable formula; the trust edges got an economy; the
City got keys. That was the plan, and it stands as V6's first strand.

\textbf{The moving ceiling.} The second strand was not planned. Between
April and June 2026, in the act of gathering, the research moved: the
reconstruction ceiling proved to be a function of time. The story of
V6's mathematics is the story of taking one letter seriously, the t that
was always in V(π, t), and following it through every term that had
quietly been treated as static.

\subsection{2. The Ceiling Moves: What R(t)
Means}\label{the-ceiling-moves-what-rt-means}

The V5.4 ceiling said: the two agents' channels, summed, cannot carry
enough information to reconstruct you. True, and proven, under two
conditions that V6 now states aloud: the agents must not collude (and
nothing may carry their residue), and the adversary must be the
adversary you measured. The second condition is where time enters. Your
archives do not change after emission. The decoders reading them do. A
model released tomorrow extracts more from yesterday's data than
anything could when the data was emitted, so the ceiling drifts upward
while you sleep, and every static guarantee has a shelf life, written
t*.

Two public events made this concrete within nine days of this document.
A four-year-old flaw in Zcash's Orchard circuit was found one day after
a frontier model's release, and weaponized the same day. A withheld
quantum optimization, attested only by a zero-knowledge proof of its
existence, was independently rediscovered in roughly two months. In both
cases the protected object sat still and the capability term moved. The
model calls this the Moving Ceiling (C82), and it is V6's headline.

\subsection{3. The Sum Leaks More Than Its Parts: Why Amnesia
Wins}\label{the-sum-leaks-more-than-its-parts-why-amnesia-wins}

In 2025 and 2026 the multi-agent privacy field proved something that
looks, at first glance, like an attack on this model: when agents pass
outputs to one another, leakage compounds, up to (2\^{}N − 1)ε for N
agents, and measurement studies found multi-agent systems leaking 68.9\%
of sensitive content overall, mostly through the unmonitored channels
between agents.

Read carefully, this is the model's own thesis in the adversary's units.
The compounding happens through the inter-agent channel. The Amnesia
Protocol's entire purpose is to delete that channel. Policy separation
asks the channel to behave and compounds exponentially; amnesia
separation removes the channel and leaks linearly. The gap between
asking and removing is, at five agents, the gap between 31ε and 5ε
(C83). The field measured the corridor this architecture bricks up.

\subsection{4. The Proof That Whispered:
Existence-Leak}\label{the-proof-that-whispered-existence-leak}

Google proved it had found a faster attack and published only the proof
of existence, hiding the method perfectly. Two months later the method
was rediscovered anyway, because the proof itself was a map with one
landmark: it said the thing exists and is findable, and every searcher
stopped digging anywhere else. V6 names this the Existence-Leak law
(C81, 70\%): a proof of feasibility leaks an upper bound on difficulty,
and there is a theorem (leakage-resilient ZK with λ \textless{} 1 is
impossible) saying the leak can never be zero. The planning corollary
(C84): every public feasibility attestation shortens your migration
deadline, whether or not the attack ever comes. Important fence: this
indicts capability attestation, not the model's own zero-knowledge
usage, which proves instances (a transaction, an attribute), not
capabilities.

\subsection{5. The Bridge and the
Forgetting}\label{the-bridge-and-the-forgetting}

Two old debts were paid in Part IV of the spec. First, the lattice and
the axes finally meet: the bridge conjecture (C85) maps the six lattice
dimensions pairwise onto the three sovereignty axes and proposes that
the recursion's base case β, the thing the two agents are allowed to
share, is exactly the First Person, which is why the separation bound
conditions on FP. The gap is β. Second, forgetting got its mathematics:
structural amnesia is conjectured to be a non-vanishing obstruction to
gluing the agents' views back into a witness (C86). Hidden things wait
for keys; forgotten things have no door. If that holds, amnesia is the
only term in the equation whose security does not erode with time, which
makes it the architectural answer to the moving ceiling of section 2.

\subsection{6. Promise Theory: The Semantic
Foundation}\label{promise-theory-the-semantic-foundation}

\subsubsection{6.1 Why Promises Matter}\label{why-promises-matter}

\textbf{Formal Spec reference:} §22

Promise Theory (Bergstra \& Burgess, 2019) provides the semantic
foundation. A promise is voluntary, unilateral, and observable.
Traditional architectures assume control: ``The server WILL do X.''
Promise architectures assume autonomy: ``The server PROMISES to do X,
and here's how you verify.''

The autonomy axiom, \emph{an agent can only make promises about its own
behaviour}, formally explains why single agents fail at the
privacy-delegation paradox. Attempting to promise both protection and
delegation exceeds autonomous capability.

\subsubsection{6.2 How Promises Map to the
Spec}\label{how-promises-map-to-the-spec}

{\def\LTcaptype{none} % do not increment counter
\begin{longtable}[]{@{}
  >{\raggedright\arraybackslash}p{(\linewidth - 2\tabcolsep) * \real{0.3333}}
  >{\raggedright\arraybackslash}p{(\linewidth - 2\tabcolsep) * \real{0.6667}}@{}}
\toprule\noalign{}
\begin{minipage}[b]{\linewidth}\raggedright
Spec Term
\end{minipage} & \begin{minipage}[b]{\linewidth}\raggedright
Promise Interpretation
\end{minipage} \\
\midrule\noalign{}
\endhead
\bottomrule\noalign{}
\endlastfoot
P (Privacy Strength) & Quality of the promise that data won't leak \\
C (Credential Verifiability) & Ability to verify a promise was kept
without seeing content \\
Φ\_agent & Separation of protection promises from delegation promises \\
Φ\_data & No single provider can break the data promise alone \\
Φ\_inference & Reasoning promises kept separate from execution
promises \\
VRC & Bilateral promise bundle between two First Persons \\
T\_∫(π) & Accumulated value of promises kept along a path \\
R(d, c, ρ) & How hard it is to break the promise retrospectively \\
The Gap (⿻) & Irreducible promise of the superagent, owned by neither
agent \\
\end{longtable}
}

V6 sharpens the last row: C78 identifies the specification-intent gap,
the distance between what an agent was told and what was meant, with the
irreducible promise of the superagent.

\textbf{Deep dive:} \texttt{promise\_theory\_reference\_v1\_4.md}

\subsection{7. Economic Architecture: VRCs and
Guilds}\label{economic-architecture-vrcs-and-guilds}

\subsubsection{7.1 Verifiable Relationship Credentials
(VRCs)}\label{verifiable-relationship-credentials-vrcs}

\textbf{Formal Spec reference:} §19 (Three Identity Layers)

A VRC is a bilateral commitment between two First Persons,
cryptographically verifiable without revealing relationship content,
relationship-scoped (dies when the relationship ends).

{\def\LTcaptype{none} % do not increment counter
\begin{longtable}[]{@{}
  >{\raggedright\arraybackslash}p{(\linewidth - 2\tabcolsep) * \real{0.5000}}
  >{\raggedright\arraybackslash}p{(\linewidth - 2\tabcolsep) * \real{0.5000}}@{}}
\toprule\noalign{}
\begin{minipage}[b]{\linewidth}\raggedright
Old Model
\end{minipage} & \begin{minipage}[b]{\linewidth}\raggedright
VRC Model
\end{minipage} \\
\midrule\noalign{}
\endhead
\bottomrule\noalign{}
\endlastfoot
Platform owns the social graph & First Persons own their edges \\
Relationships are platform assets & Relationships are bilateral
property \\
Exit means losing connections & Exit means taking your edges with you \\
\end{longtable}
}

No platform can hold relationships hostage. Switching costs collapse to
zero. Network effects accrue to people, not platforms.

\textbf{Economic parameters:}

{\def\LTcaptype{none} % do not increment counter
\begin{longtable}[]{@{}lll@{}}
\toprule\noalign{}
Parameter & Value & Purpose \\
\midrule\noalign{}
\endhead
\bottomrule\noalign{}
\endlastfoot
Ceremony & 1 ZEC (\textasciitilde\$500) & One-time genesis of agent
pair \\
Signal & 0.01 ZEC (\textasciitilde\$5) & Ongoing proof of
comprehension \\
Fee split & 61.8\% transparent / 38.2\% shielded & φ-derived constant \\
\end{longtable}
}

V6 attaches the clock here too: trust edges decay with a half-life
measured from inscription (C30 to C33), so a VRC is not a permanent
asset but a maintained one, and the signal cadence is the maintenance.
And the presence knot 🪢 carries its regime-1 declaration: it is
non-transferable, non-attesting local color, earned by walking; if it is
ever asked to attest, witnesses co-sign first.

\textbf{Deep dive:} \texttt{vrc\_promise\_protocol\_v3\_3.md}

\subsubsection{7.2 Guild Efficiency}\label{guild-efficiency}

\textbf{Formal Spec reference:} §6

A guild is a group of agents sharing a reasoning library (shared-parent
pattern). Members coordinate at O(1) cost instead of O(N²). The spec's
G(guilds) captures the multiplier.

\textbf{Trust Tier Progression:}

{\def\LTcaptype{none} % do not increment counter
\begin{longtable}[]{@{}llll@{}}
\toprule\noalign{}
Tier & Signals & Capabilities & Trust Value \\
\midrule\noalign{}
\endhead
\bottomrule\noalign{}
\endlastfoot
Blade 🗡️ & 0--50 & Basic participation, learning & 0.0--0.2 \\
Light 🛡️ & 50--150 & Multi-site coordination & 0.2--0.5 \\
Heavy ⚔️ & 150--500 & Template creation, governance & 0.5--0.8 \\
Dragon 🐉 & 500+ & Guardian eligibility, unlimited VRCs & 0.8--1.0 \\
\end{longtable}
}

\textbf{Deep dive:} \texttt{vrc\_promise\_protocol\_v3\_3.md} (§3)

\subsection{8. The Key, the Knot, and the
Star}\label{the-key-the-knot-and-the-star}

The City Key arc entered the formal lineage. The trust recursion the
City shipped in May (walk the domains, charge the trace, carry the
deepened key back) is structurally a folding scheme: the Key is an
accumulator of proofs, and the conjecture (C87) is that a real
incrementally-verifiable realization exists, with lattice-based folding
as the post-quantum hedge. The presence knot 🪢 received its honest
regime: it is local color, earned by walking, attesting nothing; if it
is ever asked to attest, witnesses co-sign first. And the eight-pointed
star gave back its borrowed gold: there is no golden ratio in that
solid, and saying so cost one paragraph and bought the geometry two
genuinely new conjectures (the parity cube, C88, and the octahedral gap,
C89: the chamber at the star's heart that neither agent owns).

\subsection{9. The Ceremonies}\label{the-ceremonies}

\subsubsection{9.1 What Ceremonies Are}\label{what-ceremonies-are}

\textbf{Formal Spec reference:} §13 (operational cycle), §15 (ceremony
implementation)

Privacy isn't just computed; it's performed. A ceremony is a structured
interaction that produces a verifiable outcome. The Celestial Ceremony
maps the operational cycle (§13) to two people:

{\def\LTcaptype{none} % do not increment counter
\begin{longtable}[]{@{}llll@{}}
\toprule\noalign{}
Phase & Symbol & Spec Operation & Human Meaning \\
\midrule\noalign{}
\endhead
\bottomrule\noalign{}
\endlastfoot
Sun & ☀️ & id(x) & Disclosure · you speak your poem \\
Gap & ⊥ & neg(x) & Silence · boundary negotiation \\
Moon & 🌑 & bnot(neg(x)) & Reflection · shared understanding \\
Return & ↻ & succ(x) & Recursion · carry forward or close \\
\end{longtable}
}

Cryptography provides guarantees. Ceremony provides meaning.

\subsubsection{9.2 The Blade Forge}\label{the-blade-forge}

\textbf{Formal Spec reference:} §15 (forge cryptography, moon phases,
tier classification)

Forging a blade: 1. Select a constellation (six sovereignty dimensions →
spectrum) 2. Walk the nodes (traverse the lattice → accumulate laps) 3.
Create the hash (SHA-256 commitment → content addressing) 4. Sign with
your key (Ed25519 binding → identity)

The spec's behavioural density ρ (§5) captures why two blades with
identical constellations but different lap counts have qualitatively
different reconstruction resistance. The Universe Blade (62 laps,
2,170s) vs Hitchhiker's Blade (13 laps, 433s): same hash position,
different proof depth.

The City Key trust recursion (C87) is the newest ceremony-grade loop:
walk the domains, charge the trace, carry the deepened Key back, each
Charge a folding step in the accumulator.

\textbf{Deep dive:} \texttt{zk\_swordsman\_blade\_forge\_v3\_0.md}

\subsection{10. Standards and the
Landscape}\label{standards-and-the-landscape}

IEEE 7012-2025 (MyTerms) published in January 2026; the implementation
effort began June 1. The EU AI Act's Annex III high-risk obligations
stand for 2026-08-02 (hedged: the Digital Omnibus would defer to
2027-12-02 but is not adopted). The model's position is unchanged and
now better armed: architecture is not against regulation, it is what
makes regulated rights enforceable. The fellow travelers (Kwaai, the
First Person Project, GliaNet, Archon, the UOR Foundation) and the
multi-agent privacy literature are cited as plurality: many teams
arriving at separation-of-duties independently, none enforcing it
architecturally, which is the model's distinctive claim and the
measurements' implicit confirmation.

\subsection{11. Conjectures in Context}\label{conjectures-in-context}

The register (\texttt{research/CONJECTURE\_REGISTER\_V6.md}) is now the
single numbering authority, born from a real failure: the corpus forked
its own numbers twice, and the V6 path's first act was to consolidate,
namespace, and promise never to renumber again. Highlights of the V6
additions: C81 Existence-Leak (70\%, held there until a second
instance), C82 the Moving Ceiling (65\%), C83 the exponential-to-linear
gap (55\%), C85 the bridge (40\%), C86 obstruction amnesia (30\%), C87
the Key accumulates (50\%). The honest-limits discipline is unchanged: λ
is unmeasured, C83's amnesia side is an engineering claim, C86 is a
framing priced at 30\%, and the model says so in its own honest-limits
section.

\subsection{12. Reading Paths by Role}\label{reading-paths-by-role}

\textbf{Reviewer or academic:} formal spec §10, §11, §5 (the conditional
relabel and R(t)), then §26 (the compounding absorption), then the
register. \textbf{Builder:} compressed spec, then §14 and §15 (amnesia
and the Key), then the regime statement before touching 🪢.
\textbf{Economist or policy reader:} this guide, then spec §27 and §30,
then the canonical figures. \textbf{City reader:} the five acts bound at
the Myth Gate, then §28, then back to Tome VIII Act 3 with the honesty
paragraph in hand. \textbf{New reader:} what-agentprivacy-is, then this
guide top to bottom, then the compressed spec.

\subsection{13. The Equation, Decoded}\label{the-equation-decoded}

Nothing in the equation's form changed at V6: the gate still multiplies,
the terms still sum, the path still integrates. What changed is the
clock attached to it. Read every term twice: once at t₀, where V5.4
proved what it proved, and once at t, where capacities drift, trust
decays from inscription, attestations discount horizons, trajectories
diverge, and exactly one term, the forgetting, holds still. The equation
was always V(π, t). V6 is the version that read the second argument.

For those who want plain English alongside the math:

\[V(\pi, t) = P^{1.5} \cdot C \cdot Q \cdot S \cdot e^{-\lambda t} \cdot (1 + A_h(\tau)) \cdot \left(1 + \sum_i w_i \frac{n_i}{N_0}\right)^k \cdot G \cdot R \cdot M \cdot \Phi_a \cdot \Phi_d \cdot \Phi_i \cdot T_{\int}(\pi)\]

{\def\LTcaptype{none} % do not increment counter
\begin{longtable}[]{@{}
  >{\raggedright\arraybackslash}p{(\linewidth - 2\tabcolsep) * \real{0.2857}}
  >{\raggedright\arraybackslash}p{(\linewidth - 2\tabcolsep) * \real{0.7143}}@{}}
\toprule\noalign{}
\begin{minipage}[b]{\linewidth}\raggedright
Term
\end{minipage} & \begin{minipage}[b]{\linewidth}\raggedright
Plain English
\end{minipage} \\
\midrule\noalign{}
\endhead
\bottomrule\noalign{}
\endlastfoot
\textbf{P\^{}1.5} & How strong is the cryptographic protection?
(superlinear · privacy compounds) \\
\textbf{C} & Can claims be verified without revealing secrets? \\
\textbf{Q} & Is the data accurate and fresh? \\
\textbf{S} & How sensitive is this data domain? \\
\textbf{e\^{}(−λt)} & How old is the data? (decays without
maintenance) \\
\textbf{(1 + A\_h(τ))} & Has this data proven itself over time?
(verified history adds value) \\
\textbf{Network term} & How connected is the sovereignty network? \\
\textbf{G} & Are agents organised into efficient guilds? \\
\textbf{R(d, c, ρ)} & How hard is it for adversaries to reconstruct?
(time-indexed as R(t) at V6) \\
\textbf{M} & How mature is the market for privacy? \\
\textbf{Φ\_agent} & Are protection and delegation properly separated?
(Swordsman ⊥ Mage) \\
\textbf{Φ\_data} & Is data distributed across providers? \\
\textbf{Φ\_inference} & Are reasoning and execution properly separated?
(Generator ⊥ Solver) \\
\textbf{T\_∫(π)} & What value accumulated along the path? (the dance,
not the stance) \\
\end{longtable}
}

\textbf{The multiplicative insight:} Any term at zero kills the whole
thing. You can't compensate for broken separation with better
cryptography. All axes must work.

\subsection{14. Quick Reference: Spec Section to Full
Context}\label{quick-reference-spec-section-to-full-context}

The V6 formal spec mirrors the V5.4 skeleton in §1 to §24, then
continues with the V6 material.

{\def\LTcaptype{none} % do not increment counter
\begin{longtable}[]{@{}
  >{\raggedright\arraybackslash}p{(\linewidth - 4\tabcolsep) * \real{0.3023}}
  >{\raggedright\arraybackslash}p{(\linewidth - 4\tabcolsep) * \real{0.1628}}
  >{\raggedright\arraybackslash}p{(\linewidth - 4\tabcolsep) * \real{0.5349}}@{}}
\toprule\noalign{}
\begin{minipage}[b]{\linewidth}\raggedright
Spec Section
\end{minipage} & \begin{minipage}[b]{\linewidth}\raggedright
Topic
\end{minipage} & \begin{minipage}[b]{\linewidth}\raggedright
Full Context
\end{minipage} \\
\midrule\noalign{}
\endhead
\bottomrule\noalign{}
\endlastfoot
§1 & The equation & this guide §13 ·
\texttt{privacy\_is\_value\_v5.md} \\
§4 & Three-axis separation &
\texttt{VISUAL\_ARCHITECTURE\_GUIDE\_v2\_0.md} \\
§5 & Reconstruction difficulty · R(t) and t* & this guide §2 ·
\texttt{dualprivacy\_researchpaper\_v4\_0.md} \\
§6 & Guild efficiency & this guide §7 ·
\texttt{vrc\_promise\_protocol\_v3\_3.md} §3 \\
§7 & Path integral T\_∫ & V5.2 + V5.3 Research Notes \\
§8 & Holographic bound &
\texttt{uor\_tetrahedra\_zk\_mapping\_v2\_0.md} \\
§10 & Separation bound · preconditions & this guide §2 ·
\texttt{dualprivacy\_researchpaper\_v4\_0.md} §5 \\
§11 & Reconstruction ceiling · conditional regime & this guide §2, §3 \\
§12 & Z/(2⁶)Z algebra & \texttt{SYSTEMS\_HEXAGRAM\_PHYSICS.md} \\
§13 & Operational cycle & this guide §9 · V5.3 Research Note \\
§14 & Amnesia Protocol · obstruction upgrade & this guide §3, §5 \\
§15 & Ceremony + Forge · the Key as accumulator & this guide §8, §9 ·
\texttt{ceremonies/} directory \\
§19 & Identity layers & \texttt{swordsman\_mage\_whitepaper\_v6\_0.md}
§7 \\
§20 & Cosmological quaternion & Grimoire v10.0.0, Act XXXI \\
§22 & Promise Theory & this guide §6 ·
\texttt{promise\_theory\_reference\_v1\_4.md} \\
§25 & The two worked instances (Orchard · Schrottenloher) & this guide
§2 \\
§26 & Compounding bound · exponential-to-linear gap & this guide §3 \\
§27 & The temporal thread & this guide §2, §4 \\
§28 & The geometry of the gap & this guide §8 \\
§29 & Narrative corpus & \texttt{tomes/} · the bound collection \\
§30 & Canonical figures & this guide §12 (economist path) \\
§31 & External landscape & this guide §10 \\
§32 & Honest limits & this guide §11 \\
§33 & References & \texttt{privacy\_value\_v6.md} §33 \\
\end{longtable}
}

\subsection{15. Glossary Bridge}\label{glossary-bridge}

Key terms that appear in the spec without full definition:

{\def\LTcaptype{none} % do not increment counter
\begin{longtable}[]{@{}
  >{\raggedright\arraybackslash}p{(\linewidth - 4\tabcolsep) * \real{0.1714}}
  >{\raggedright\arraybackslash}p{(\linewidth - 4\tabcolsep) * \real{0.3429}}
  >{\raggedright\arraybackslash}p{(\linewidth - 4\tabcolsep) * \real{0.4857}}@{}}
\toprule\noalign{}
\begin{minipage}[b]{\linewidth}\raggedright
Term
\end{minipage} & \begin{minipage}[b]{\linewidth}\raggedright
Spec Usage
\end{minipage} & \begin{minipage}[b]{\linewidth}\raggedright
Full Definition
\end{minipage} \\
\midrule\noalign{}
\endhead
\bottomrule\noalign{}
\endlastfoot
\textbf{First Person} & Implied throughout & The human whose behavioural
capital is at stake. Not ``user''; users are used. \\
\textbf{Sovereignty} & ``sovereignty lattice'' & Self-determination over
one's own boundaries, delegations, and data \\
\textbf{Blade} & §12, §15 & A forged commitment, six dimensions,
cryptographically bound \\
\textbf{Constellation} & §15 & The selection of dimensions before
forging \\
\textbf{Stratum} & §12, §15 & Hamming weight: how many sovereignty
dimensions are active \\
\textbf{Spectrum} & §12 & Six-bit decomposition: which specific
dimensions \\
\textbf{Datum} & §12 & Raw blade value (0--63) \\
\textbf{The Gap} (⿻) & §10 & The irreducible space where sovereignty
lives, not a void but a guarantee \\
\end{longtable}
}

V6 terms:

{\def\LTcaptype{none} % do not increment counter
\begin{longtable}[]{@{}
  >{\raggedright\arraybackslash}p{(\linewidth - 4\tabcolsep) * \real{0.1714}}
  >{\raggedright\arraybackslash}p{(\linewidth - 4\tabcolsep) * \real{0.3429}}
  >{\raggedright\arraybackslash}p{(\linewidth - 4\tabcolsep) * \real{0.4857}}@{}}
\toprule\noalign{}
\begin{minipage}[b]{\linewidth}\raggedright
Term
\end{minipage} & \begin{minipage}[b]{\linewidth}\raggedright
Spec Usage
\end{minipage} & \begin{minipage}[b]{\linewidth}\raggedright
Full Definition
\end{minipage} \\
\midrule\noalign{}
\endhead
\bottomrule\noalign{}
\endlastfoot
\textbf{Moving ceiling} & §5, C82 & R(t) is non-decreasing under
capability growth; the ceiling drifts while the archives sit still \\
\textbf{Shelf life t}* & §5 & sup\{t : R(t) \textless{} 1\}: the last
moment a static guarantee still holds \\
\textbf{Existence-leak} & §27, C81 & A public proof of feasibility leaks
an upper bound on difficulty; the leak can never be zero \\
\textbf{Ages progressively} & §27, C47 & Taxonomy class for guarantees
that weaken continuously as decoders improve, rather than failing at one
event \\
\textbf{Obstruction amnesia} & §14, C86 & Grade-2 forgetting as a
non-vanishing obstruction class to gluing the agents' local views into a
global witness \\
\textbf{Terminal obstruction} & §14, C73 & Structural loss of β, the
recursion's base case; no successor state exists \\
\textbf{Parity cube} & §28, C88 & The stella octangula's two tetrahedra
as the even/odd parity classes of \{0,1\}³ \\
\textbf{Octahedral gap} & §28, C89 & The intersection core as the locus
of the conditional-independence bound; the chamber neither agent owns \\
\textbf{The Key accumulates} & §15, C87 & The City Key trust recursion
as incrementally verifiable computation; each Charge is a folding
step \\
\textbf{Regime 1} & §15 & The presence knot's declared status:
non-transferable, non-attesting local color, earned by walking \\
\textbf{The register} & throughout &
\texttt{CONJECTURE\_REGISTER\_V6.md}, the single numbering authority,
head C89 \\
\end{longtable}
}

\textbf{Full glossary:} \texttt{GLOSSARY\_MASTER\_v4\_0.md} (160+
entries)

\subsection{Document Metadata}\label{document-metadata}

{\def\LTcaptype{none} % do not increment counter
\begin{longtable}[]{@{}
  >{\raggedright\arraybackslash}p{(\linewidth - 2\tabcolsep) * \real{0.5000}}
  >{\raggedright\arraybackslash}p{(\linewidth - 2\tabcolsep) * \real{0.5000}}@{}}
\toprule\noalign{}
\begin{minipage}[b]{\linewidth}\raggedright
Field
\end{minipage} & \begin{minipage}[b]{\linewidth}\raggedright
Value
\end{minipage} \\
\midrule\noalign{}
\endhead
\bottomrule\noalign{}
\endlastfoot
Version & 6.0 (unified V6 labeling across all canon papers, G3 decision
2026-06-10) \\
Status & Fully standalone: V5.4 companion material (Promise Theory,
economics, ceremonies, quick reference, glossary) carried in full and
updated to V6 \\
Authority & privacy\_value\_v6.md · CONJECTURE\_REGISTER\_V6.md (head
C89) \\
Siblings & pvm\_v6\_compressed.md (Swordsman) · research paper V6
edition · whitepaper V6 edition \\
Working record & research/privacy\_value\_v6\_draft.md (Parts I to V,
full honest limits) \\
License & CC BY-SA 4.0 \\
\end{longtable}
}

(⚔️⊥⿻⊥🧙)😊

\end{document}
